% Le CV de Joel Codja, en français. Construit par cv/build.sh dans public/cv/joel-codja-cv-fr.pdf.
% Une page, une colonne, des titres standard : les logiciels de recrutement le lisent comme du texte.
\documentclass[10pt]{article}

% Shared layout of both resumes (cv/joel-codja-cv.tex, cv/joel-codja-cv-fr.tex).
\usepackage[T1]{fontenc}
\usepackage[utf8]{inputenc}
\usepackage{charter}
\usepackage[a4paper,margin=1.6cm,top=1.4cm,bottom=1.4cm]{geometry}
\usepackage{titlesec}
\usepackage{enumitem}
\usepackage{tabularx}
\usepackage{xcolor}
\usepackage[hidelinks]{hyperref}

% Real text in the PDF, so applicant tracking systems can read and copy it.
\input{glyphtounicode}
\pdfgentounicode=1

% The site's terracotta, darkened to print well on white.
\definecolor{accent}{HTML}{B4532F}
\definecolor{muted}{HTML}{5A5A5A}

\pagestyle{empty}
\setlength{\parindent}{0pt}
\setlength{\parskip}{0pt}
\linespread{1.06}
\raggedright

% Section: a terracotta label over a hairline (plain case: small caps split the text for ATS).
\titleformat{\section}{\large\bfseries\color{accent}}{}{0pt}{}[\vspace{-6pt}\color{muted!40}\rule{\linewidth}{0.4pt}]
\titlespacing*{\section}{0pt}{11pt}{6pt}
\newcommand{\cvsection}[1]{\section*{#1}}

\setlist[itemize]{leftmargin=1.1em,itemsep=1.5pt,topsep=3pt,parsep=0pt,label={\color{muted}\textbullet}}
\newenvironment{cvitems}{\begin{itemize}}{\end{itemize}}

% Name, title, then one line of contacts as plain text links: {title}{location}
\newcommand{\cvheader}[2]{%
  {\Huge Joel Codja}\par\vspace{3pt}
  {\large\color{accent}#1}\par\vspace{5pt}
  {\small\color{muted}#2 \,\textbar\, +229 01 63 04 15 88 \,\textbar\,
  \href{mailto:gide749@gmail.com}{gide749@gmail.com}\\
  \href{https://joelcodja.com}{joelcodja.com} \,\textbar\,
  \href{https://github.com/cjdotdev}{github.com/cjdotdev} \,\textbar\,
  \href{https://www.linkedin.com/in/joel-codja-4356a2361/}{linkedin.com/in/joel-codja-4356a2361}}\par
}

% {company}{link}{role}{dates}
\newcommand{\cventry}[4]{%
  {\bfseries #1} \,{\small\color{muted}#2}\hfill{\small #4}\par
  {\itshape #3}\par\vspace{4pt}
}
% {name}{link}
\newcommand{\cvproject}[2]{%
  \vspace{2pt}{\bfseries #1} \,{\small\color{muted}#2}\par\vspace{1pt}
}
\newcommand{\cvstack}[1]{\vspace{1pt}{\small\color{muted}#1}\par\vspace{5pt}}

\newenvironment{cvskills}{\tabularx{\linewidth}{@{}>{\bfseries}p{3.6cm}X@{}}}{\endtabularx}


\begin{document}

\cvheader{Développeur backend}{Cotonou, Bénin}

\cvsection{Profil}
Développeur backend, orienté conception d'API et architectures serverless edge.
Ouvert à un poste backend à temps plein comme aux missions en freelance.

\cvsection{Expérience}
\cventry{Lixali Studio}{\href{https://lixalistudio.tech}{lixalistudio.tech}}{Développeur backend}{Avr. 2026 -- Aujourd'hui}
\cvproject{MakiSend}{\href{https://github.com/lixalistudio-oss/makisend}{github.com/lixalistudio-oss/makisend}}
API open source d'emails transactionnels qui tourne sur le compte Cloudflare du client et envoie via
son propre fournisseur SMTP. Développeur backend principal (156 commits sur 224).
\begin{cvitems}
  \item Conception du pipeline d'envoi asynchrone : l'API répond \texttt{202} tout de suite, puis un
        Workflow Cloudflare durable livre en SMTP, avec 3 nouvelles tentatives à 30\,s d'intervalle et
        une dead-letter queue. 300 à 400\,ms par envoi au mieux.
  \item Chiffrement des identifiants, destinataires, contenus et pièces jointes en AES-256-GCM avant
        leur arrivée en base, et rotation des clés qui rechiffre les valeurs stockées par
        identifiant de clé.
  \item Construction de l'API REST avec Hono : validation Zod, clés d'idempotence, quotas d'envoi par clé
        et une référence OpenAPI 3.1 générée depuis les mêmes schémas Zod.
  \item Plus de 400 tests Vitest unitaires, d'intégration et end-to-end, dont un
        vrai serveur SMTP, lancés par GitHub Actions à chaque pull request.
\end{cvitems}
\cvstack{TypeScript, Hono, Cloudflare Workers, D1, Queues, Workflows, Drizzle, Zod, Better Auth, Vitest}

\cvsection{Projets}
\cvproject{linkedin-scheduler}{\href{https://github.com/cjdotdev/linkedin-scheduler}{github.com/cjdotdev/linkedin-scheduler}}
Outil open source qui publie des posts LinkedIn écrits à l'avance en fichiers Markdown, via l'API
officielle de LinkedIn. Ni serveur ni interface web : un dossier, deux scripts et une ligne cron.
\begin{cvitems}
  \item Ne publie jamais deux fois : un registre, un état de publication et un verrou entre les
        exécutions ; \texttt{-{}-check} vérifie le jeton et \texttt{-{}-dry-run} montre ce qui partirait.
  \item Génère des visuels 1080$\times$1350 depuis un template HTML, avec des extraits de code colorés.
\end{cvitems}
\cvstack{Python, cron, API LinkedIn, Node.js, Playwright, Shiki}

\cvproject{Articles techniques}{\href{https://joelcodja.com/fr/blog/}{joelcodja.com/fr/blog}}
Articles en français et en anglais : le chiffrement au repos AES-256-GCM en Node.js, la validation
des entrées d'une API avec Zod, Drizzle ORM pour les backends TypeScript.

\cvsection{Compétences}
\begin{cvskills}
  Langages & TypeScript, SQL, Python \\
  API & Hono, Node.js, REST, OpenAPI, Scalar, Zod \\
  Cloud et données & Cloudflare Workers, D1, Queues, Workflows, Drizzle, SQLite, PostgreSQL \\
  Qualité et livraison & Vitest, GitHub Actions, Docker \\
\end{cvskills}

\cvsection{Langues}
Français (langue maternelle), anglais (C1)

\end{document}
